% --- Back cover ---
% Large PERSONIX wordmark, no logo image. Cream paper inherited globally.
% Absolute (tikz overlay) positioning so the footer sticks to the bottom.
\clearpage
\thispagestyle{empty}
\begin{tikzpicture}[remember picture, overlay]
  % --- Wordmark ---
  \node[anchor=north, font=\sffamily\bfseries\fontsize{48}{52}\selectfont]
    at ([yshift=-26mm]current page.north) {PERSONIX};
  \node[anchor=north, font=\rmfamily\itshape\large]
    at ([yshift=-44mm]current page.north) {Changement sans compromis};
  \node[anchor=north, font=\rmfamily\small,
        text width=\paperwidth-50mm, align=center]
    at ([yshift=-54mm]current page.north)
    {Un réseau de réputation décentralisé, incensurable et incorruptible};
  % --- Body ---
  \node[anchor=north, text width=\paperwidth-50mm, align=justify, font=\small]
    at ([yshift=-72mm]current page.north)
    {%
      L'État fonctionnait quand la communication était lente, les décisions
      locales, et l'autorité vivait dans la même ville que celle qu'elle
      gouvernait. Les réseaux d'aujourd'hui inversent ces trois hypothèses
      --- et les institutions bâties sur elles ne conviennent plus au monde
      qu'elles régissent. Ce livre décrit un outil qui, lui, convient :
      un réseau de réputation décentralisé où l'identité est vôtre à garder,
      où la vérité est publiquement auditable, et où la corruption est un
      suicide réputationnel.

      \smallskip
      Pas un manifeste. Pas une prévision. Un plan opérationnel de la manière
      dont le successeur de l'État peut être bâti --- volontairement, une
      déclaration à la fois.
    };
  % --- Footer (anchored to the bottom of the page) ---
  \node[anchor=south, font=\sffamily\footnotesize,
        text width=\paperwidth-50mm, align=center]
    at ([yshift=12mm]current page.south)
    {%
      Pavel Kudrna\quad\textbullet\quad Prague, 2026\par
      \href{https://personix.org}{personix.org}\par
      Sous licence Creative Commons Attribution-ShareAlike 4.0 International
      (CC BY-SA 4.0).
    };
\end{tikzpicture}%
\mbox{}%
\clearpage
